% State of the art --- section structure from prd.md 3.1 (Ingenieur, Chapter 3), ownership from
% Table 3. Rebuilt 26 Sep in the Master's form: each section is a taxonomy, ONE comparison table
% whose last row is "This work", and a "Critical comparison." paragraph that compares the studies
% TO EACH OTHER. The studies are compared to this work only in sec:research-gaps, which is why the
% positioning against Lim et al. and the concessions to CommandSwarm, SkySim and MIRA live there.
%
% This chapter OWNS: the voice-UAV, language-to-robot and dual-path literature (this document
% only), the positioning against Lim et al., and the research gaps.
% It does NOT own: edge inference, quantisation, constrained decoding or spoken-language
% understanding (Master, State of the art) -- one sentence and a cross-reference at most.
% The choice of each speech component against this system's requirements is a design decision and
% lives in chap:architecture (sec:component-choices), not here.
% A dash in a table marks what the review did not establish from the source; a wrong cell is a
% BLOCKER. Every cell restates a fact the chapter prose (reviewed in issues 07 and 11) already
% carried, or a fact recorded for its key in .scratch/thesis-bibliography/ingenieur-candidates.md.

\chapter{State of the art}
\label{chap:state-of-the-art}
This chapter reviews the four bodies of published work the system of
Chapter~\ref{chap:introduction} resembles: language interfaces to aerial and other robots,
architectures that pair a fast path with a slower one, the control of vehicle swarms, and the speech
front ends of spoken-command systems. Each section first sorts its studies into the kinds that
exist, then compares them to one another in a table along fixed axes, and closes by reading that
table. The reading states what the studies agree on, where they differ, and what none of them
provides. Section~\ref{sec:research-gaps} combines the four readings into three numbered gaps, maps
each onto the contribution or the design decision that closes it, and positions this work against
its baseline. The studies are compared to each other in the tables and to this work only in that
final section. The foundations the tables assume are those of Chapter~\ref{chap:background}. The
literature on edge inference, quantisation, constrained decoding and spoken-language understanding
is reviewed in the \emph{M\'emoire de Master} and is not repeated here.

\section{Voice-controlled UAV systems}
\label{sec:voice-uav}
Language interfaces to \glspl{uav} have moved from recognisers matched against a fixed command
language to \glspl{llm} that plan, generate programs or call flight functions, and a recent survey
maps that field and the directions it is taking~\cite{javaid2024}. The published systems are of four
kinds: recognisers matched against a fixed command language, general language-model control of
robots, agents built for a single drone, and interfaces to swarms. Four properties separate them:
whether the input is speech or typed text, whether one vehicle or a group is commanded, where the
language model runs, and what check, if any, stands between the model's output and the vehicle.
Table~\ref{tab:voice-uav-survey} records the four properties for the systems reviewed in this
section.

\paragraph{Fixed command languages.} The earlier generation of voice control needed no language
model, because it confined the operator to a command language fixed in advance. Landau and van
Delden map the output of the Nuance speech-recognition platform onto flight commands for a DJI
Phantom~4 through regular expressions that define the control language~\cite{landau2017}.
Contreras et al.\ improve a cloud-based recogniser for \gls{uav} control by matching its output
against the phonemes of a domain vocabulary, which raises English command recognition from 74.81\%
to 93.33\% in a simulated domestic scene. They evaluate it with several levels of distortion applied
to the spoken commands~\cite{contreras2020}. Both systems command a single vehicle.

\paragraph{Language-model control.} \Glspl{llm} replaced the fixed command language with generated
plans, programs or function calls, and the check between the model's output and actuation became
the central design question. SayCan constrains a language model to the actions a robot can perform
by weighting its proposals with the learned value functions of pretrained skills~\cite{saycan2022}.
Code as Policies generates programs that compose given perception and control primitives. Its
authors state that the approach assumes every given instruction is feasible and cannot tell in
advance whether a response will be correct~\cite{codeaspolicies2023}. Vemprala et al.\ apply
ChatGPT, through prompt design and a library of robot functions, to tasks that include aerial
navigation. The check there is a human on the loop, who evaluates the output for quality and safety
before it is executed, and the authors advise that such tools should not be given full control of
the robotics pipeline~\cite{chatgptrobotics2024}.

\paragraph{Single-vehicle drone agents.} Systems built specifically for drones make the latency of
generation explicit. TypeFly observes that sequential token generation introduces a latency that
grows with the length of the plan, and it answers by changing the output representation. A
token-efficient plan language, interpreted while it streams, reduces response time by up to 62\% and
keeps it below 1.5~s on a benchmark of eleven tasks, with GPT-4 as a remote model and a typed task
description as input~\cite{typefly2024}. Chat with UAV splits planning from execution between two
language-model agents~\cite{chatwithuav2025}. Phadke et al.\ couple existing language models to
robotic simulators through a modular template framework~\cite{phadke2024}. Silva and Burke report
what they describe as the broadest evaluation of a language-model command interface for drones on
the axes it measures. It covers more than 1{,}000 simulated flights, 110 geofence-violation trials
in which no aircraft left the permitted zone, and three real quadcopters, with the language model
treated as an untrusted commander whose output is validated on the server~\cite{silva2026}. The
models that evaluation uses are mostly those of cloud providers. The agent of Lim et
al.~\cite{lim2025} prompts open models, Gemma3, Qwen2.5, Llama-3.2 and DeepSeek-LLM, zero-shot
through Ollama on a local workstation \gls{gpu}, and flies one PX4-based quadcopter to a goal found
by a vision-language model; Section~\ref{sec:research-gaps} examines it as the baseline of this
document.

\paragraph{Where the model runs.} Two systems make the host of the language model their subject.
Sikorski et al.\ command a wheeled mobile robot by speech, recognised offline with VOSK. They compare
GPT-4-Turbo in the cloud against LLaMA~2-7B quantised to Q5\_K\_M and run offline on a control
computer, and the offline model exhibits significant limitations in consistency and
reliability~\cite{sikorski2025}. Torkamani and Zarin route spoken commands for IoT devices between
edge and cloud according to processor load, device temperature and network latency, with the edge
tier on an NVIDIA Jetson platform~\cite{torkamani2025}.

\paragraph{Swarms.} Language interfaces to swarms are more recent. SwarmChat offers voice, text and
teleoperation interfaces to a robotic swarm over ROS~2 through modules described as based on
language models. In its implementation, intent recognition is rule-based keyword matching, which its
authors list as a limitation, and its evaluation is preliminary~\cite{swarmchat2025}. Iannoli et
al.\ evaluate six general-purpose language models on four swarm missions from natural-language
text, in ArduPilot software-in-the-loop simulation. They find that the models still struggle to
execute even simple swarm tasks reliably without explicit grounding and execution support, and that
task-specific planning tools and runtime guardrails substantially improve
robustness~\cite{iannoli2026}. SkySim plans waypoints for a swarm with a cloud-hosted model above a
potential-field layer~\cite{skysim2026}, and CommandSwarm turns speech or text into behaviour trees
for a robotic swarm, gated by a deterministic parser~\cite{commandswarm2026};
Section~\ref{sec:swarm-control-sota} compares their control layers.

\begin{table}[htbp]
\centering
\footnotesize
\caption[Published language interfaces to robots]{Published language interfaces to robots, by
the four properties of Section~\ref{sec:voice-uav}. \emph{Remote}: a hosted model reached over a
network. \emph{Check before execution}: the mechanism between the model's (or the recogniser's)
output and the vehicle, as the source describes it, or, where the source describes none, the
constraint it places on generation. A dash marks a property this review did not establish from the
source. The last row is this work.}
\label{tab:voice-uav-survey}
\begin{tabularx}{\textwidth}{@{}>{\raggedright\arraybackslash}p{2.5cm} >{\raggedright\arraybackslash}p{1.7cm} >{\raggedright\arraybackslash}p{2.0cm} >{\raggedright\arraybackslash}X >{\raggedright\arraybackslash}X@{}}
\toprule
System & Input & Target & Language path and host & Check before execution \\
\midrule
Landau and van Delden~\cite{landau2017} & Speech & One \gls{uav} & Nuance recogniser; no language model & Regular-expression command language \\
Contreras et al.~\cite{contreras2020} & Speech & One simulated \gls{uav} & Cloud recogniser; no language model & Matching against domain phonemes \\
Vemprala et al.~\cite{chatgptrobotics2024} & Text & Robots, including aerial & ChatGPT, remote & Human on the loop \\
TypeFly~\cite{typefly2024} & Text & One \gls{uav} & GPT-4, remote & Plan language; no check stated \\
Silva and Burke~\cite{silva2026} & --- & \glspl{uav} & Mostly cloud models & Server-side validation; geofence \\
Lim et al.~\cite{lim2025} & Text & One \gls{uav} & Open models, local \gls{gpu} workstation & Prompt format; temperature 0.2 \\
Sikorski et al.~\cite{sikorski2025} & Speech & One ground robot & GPT-4-Turbo, remote, or LLaMA~2-7B on a control computer & --- \\
Torkamani and Zarin~\cite{torkamani2025} & Speech & IoT devices & Edge (Jetson) or cloud, routed at run time & --- \\
SwarmChat~\cite{swarmchat2025} & Speech, text, teleoperation & Robot swarm & Modules based on language models; --- & Rule-based intent keywords \\
Iannoli et al.~\cite{iannoli2026} & Text & \gls{uav} swarm & Six general-purpose models; --- & Planning tools; runtime guardrails \\
SkySim~\cite{skysim2026} & Text & \gls{uav} swarm, 3--30 & Gemini 3.5 Pro, remote & Waypoint bounds; potential-field layer \\
Command\-Swarm~\cite{commandswarm2026} & Speech or text & Robot swarm & Open models of 6.7--14B, 4-bit; --- & Primitive whitelist; parser; safety classifier \\
\midrule
This work & Speech & Five simulated \glspl{uav} & Fine-tuned 0.5B model; Raspberry~Pi~5 \gls{cpu} & Grammar; validator; flight state machine \\
\bottomrule
\end{tabularx}
\end{table}

\paragraph{Critical comparison.} Read down its columns, Table~\ref{tab:voice-uav-survey} exhibits
four patterns. On input, speech and typed text divide the table roughly in half, and speech input
occurs with a group of vehicles in only two systems, SwarmChat and CommandSwarm. Neither states the
hardware on which its model runs. On the target, four systems address a swarm, and only SkySim's row
records the number of vehicles flown. On the host, the
language model is remote in every system that names a provider, local to a workstation \gls{gpu} in
Lim et al., on a control computer in the offline arm of Sikorski et al., and on an edge
accelerator in Torkamani and Zarin, which hands work to the cloud when its runtime metrics call for
it. No system in the table runs its full chain on an \gls{sbc} without a \gls{gpu}. On the check,
the systems without a language model constrain the operator instead, through a command language or a
phoneme vocabulary. Every system with a language model constrains or checks its output in some way.
The checks range from a prompt format and a sampling temperature (Lim et al.) and human review
(Vemprala et al.) to server-side validation (Silva and Burke) and a whitelist parser behind a safety
classifier (CommandSwarm). The checks that are deterministic and automated appear in the systems of
2026, and Iannoli et al.\ measure why: without planning tools and guardrails, general-purpose
models do not execute swarm tasks reliably. None of the systems in the table reports a stopping
command handled separately from the path that interprets commands.

\section{Fast-path and dual-path architectures}
\label{sec:soa-dual-path}
Several systems split the handling of speech between a fast path and a slower, more capable one, and
they are of three kinds. In the first, the two paths are \emph{alternatives}: a study builds a direct
pipeline and a cascade for the same task, measures both, and selects one. In the second, the fast
path is a cheaper branch \emph{inside one model}, selected frame by frame. In the third, the two paths
run \emph{concurrently}, and a gate or a preemption decides which one's output takes effect. A fourth
arrangement divides the work between two components by function rather than by speed.

\paragraph{Alternatives compared.} Two studies measure the latency cost of the pipeline's shape
directly. Sim\~oes et al.\ compare, for one Tello drone commanded in Portuguese, a cascade of speech
recognition and a language model against a classifier that maps speech directly to a command. The
cascade reaches an accuracy of 0.81 at 1.233~s, the direct classifier 0.99 at
0.021~s~\cite{simoes2024}. Henry et al.\ report, on the explicit simple commands of a French corpus
of spontaneous human--drone interaction, 93\% at 7~ms for a transcription-free intent recogniser
against 79\% at 202~ms for a cascade of Whisper and CamemBERT-Large. They attribute that 14-point gap
to the cascade architecture rather than to recognition quality. On the full spontaneous test set the
figures are 82\% and 59\%, and the 7~ms is measured on a \gls{gpu} (106~ms on a
\gls{cpu})~\cite{henry2026}.

\paragraph{A branch inside one model.} Sahai et al.\ use wake-word detection to select, frame by
frame, which branch of the attention networks inside one recognition model is executed. The
wake-word frames receive the cheaper branch, in order to save computation~\cite{dualattn2023}.

\paragraph{Concurrent paths.} RelayS2S pairs a fast duplex speech-to-speech model, which drafts the
opening of a response, with a slower cascade of recognition and a language model. A lightweight
learned verifier decides whether the fast draft is committed, and the 90th-percentile latency of the
first chunk falls from 1{,}006~ms to 81~ms~\cite{relays2s}. MIRA combines low-latency barge-in
preemption and streaming response generation with deliberative turn decisions in a dual-timescale
interaction policy~\cite{mira}. It bounds the robot's physical commitment to a short, interruptible
prefix and enforces physical safety constraints in a robot-side layer at the control rate. Its fast
interruption gate works from voice activity alone, without transcription. When the user speaks
during the robot's playback, it dispatches an abort, stops the speech output and ends the robot's
motion, after which the robot holds its last safe pose. MIRA addresses dialogue with a single
embodied companion, and it does not name its recognition and language models or say where they
run.

\paragraph{A split by function.} Chat with UAV divides planning and execution between two
language-model agents~\cite{chatwithuav2025}. The division is by role, and both of its branches run
at the latency of a language model. Table~\ref{tab:soa-dual-path} sets the four arrangements side by
side.

\begin{table}[htbp]
\centering
\footnotesize
\caption[Architectures with a fast and a slow path]{Architectures that pair a fast path with a
slower one, by how the paths are arranged and what the fast one decides. \emph{Alternatives}: both
paths are built and measured, and one is selected. \emph{Concurrent}: both run on the same input,
and a gate or a preemption decides which output takes effect. Latencies are those the source
reports for the fast and the slow path, in that order. A dash marks a property this review did not
establish from the source. The last row is this work.}
\label{tab:soa-dual-path}
\begin{tabularx}{\textwidth}{@{}>{\raggedright\arraybackslash}p{1.9cm} >{\raggedright\arraybackslash}X >{\raggedright\arraybackslash}p{2.0cm} >{\raggedright\arraybackslash}X >{\raggedright\arraybackslash}p{2.35cm} >{\raggedright\arraybackslash}p{1.9cm}@{}}
\toprule
Work & Fast path; slow path & Arrangement & What the fast path decides & Purpose & Latency reported \\
\midrule
Sim\~oes et al.~\cite{simoes2024} & Direct speech-to-command classifier; recognition then a language model & Alternatives & The command & --- & 0.021~s; 1.233~s \\
Henry et al.~\cite{henry2026} & Transcription-free intent recogniser; Whisper then CamemBERT-Large & Alternatives & The intent & --- & 7~ms (\gls{gpu}), 106~ms (\gls{cpu}); 202~ms \\
Sahai et al.~\cite{dualattn2023} & Cheaper attention branch; full branch, in one recognition model & One model, per frame & Which branch runs on a frame & Computation & --- \\
RelayS2S~\cite{relays2s} & Duplex speech-to-speech draft; recognition then a language model & Concurrent; learned verifier & The opening of a response, if committed & Responsiveness & 81~ms; 1{,}006~ms at p90, first chunk \\
MIRA~\cite{mira} & Barge-in gate on voice activity; deliberative turn decisions & Concurrent; preemption & One action: abort the response and stop motion & Barge-in & --- \\
Chat with UAV~\cite{chatwithuav2025} & None: planning agent and execution agent, both language models & Split by function & --- & --- & --- \\
\midrule
This work & Two-class keyword spotter; endpointer, recognition, language model under a grammar, validator & Concurrent; preemption by sequence number & One of two swarm commands, a hold or an abort & Stopping the swarm & 150~ms; 2{,}500~ms, p95 targets \\
\bottomrule
\end{tabularx}
\end{table}

\paragraph{Critical comparison.} Table~\ref{tab:soa-dual-path} separates the works along four axes.
On the arrangement, the two drone studies build their paths as alternatives and select between them,
so neither has to decide what happens when both paths are running. The within-model branch and the
concurrent systems do decide it: by a per-frame switch in one model, a learned verifier, or a
preemption. On what the fast path decides, the rows grow narrower down the table. The alternative
pipelines decide a full command or intent; the verifier of RelayS2S commits a drafted response;
MIRA's gate has a single action. On purpose, the fast paths of Sahai et al.\ and RelayS2S serve computation
and responsiveness. MIRA's gate stops speech and motion on barge-in, and MIRA assigns its physical
safety constraints to a separate robot-side layer at the control rate. On latency, the two
drone studies agree that removing transcription cuts the latency of a command by one to two orders
of magnitude, 21~ms against 1{,}233~ms and 7~ms against 202~ms. Henry et al.\ also report the direct
path as the more accurate of the two, by 14 points on explicit simple commands, so on these tasks the
faster shape costs no accuracy. That result bears on a fast path that must choose among several
commands, and it is reported for alternatives, never for two paths running at once. No work in the
table states a rule deciding which commands a fast path may carry.

\section{Swarm control}
\label{sec:swarm-control-sota}
The control of a swarm draws on three kinds of work: coordination models that produce group motion
from rules applied by each vehicle, separation methods that keep the vehicles apart, and, more
recently, architectures that place a language model above a swarm controller. The first two are
foundations, set out in Section~\ref{sec:bg-swarm}: Reynolds' flocking rules~\cite{reynolds1987} and
the artificial potential field, whose failure modes Koren and Borenstein identify as inherent to the
method~\cite{koren1991}. What this section compares is how the third kind combines a language model
with them.

\paragraph{Integration of a language model.} Strobel et al.\ distinguish two ways of integrating a
language model with a swarm. In \emph{indirect} integration, a language model synthesises and
validates controllers before or during deployment. In \emph{direct} integration, a separate
language-model instance runs on each robot. They describe their contribution as mainly
conceptual~\cite{llm2swarm2024}. SkySim places a cloud-hosted planner, Gemini 3.5 Pro, above an
artificial-potential-field layer running at 20~Hz that repels vehicles from one another, with speed
saturated at 0.5~m/s, for swarms of 3, 10 and 30 simulated drones in Gazebo with
ROS~2~\cite{skysim2026}. Geofencing is performed by the planner, which rejects waypoints outside the
permitted zone and falls back to holding position. A malformed output or a timeout of the model's
interface instead keeps the waypoints of the previous command. The mean planning latency rises from
34~s for 3 drones to 50~s for 10, with outliers above 100~s, a latency its authors judge unsuitable
for time-critical missions. CommandSwarm generates behaviour trees from a whitelist of executable
swarm primitives. It targets a PyGame-based swarm simulator but scores its generated trees by
text-similarity metrics and parser acceptance rather than executing them~\cite{commandswarm2026}.
Iannoli et al.\ run general-purpose models in ArduPilot software-in-the-loop simulation with
planning tools and runtime guardrails~\cite{iannoli2026}. Table~\ref{tab:soa-swarm} compares the
control layers of these works.

\begin{table}[htbp]
\centering
\footnotesize
\caption[Swarm control beneath a language interface]{Coordination models and swarm architectures
with a language interface, by what coordinates the vehicles, what keeps them apart, and where the
language model sits. \emph{Evaluation}: how the control layer was exercised. A dash marks a
property this review did not establish from the source. The last row is this work.}
\label{tab:soa-swarm}
\begin{tabularx}{\textwidth}{@{}>{\raggedright\arraybackslash}p{2.1cm} >{\raggedright\arraybackslash}X >{\raggedright\arraybackslash}p{2.4cm} >{\raggedright\arraybackslash}X >{\raggedright\arraybackslash}X@{}}
\toprule
Work & Coordination & Separation & Role of the language model & Evaluation \\
\midrule
Reynolds~\cite{reynolds1987} & Three local rules: separation, velocity matching, centring & Collision-avoidance rule & None & --- \\
Koren and Borenstein~\cite{koren1991} & --- & Potential field; failure modes analysed & None & --- \\
Strobel et al.~\cite{llm2swarm2024} & --- & --- & Indirect: synthesises and validates controllers; or direct: one instance per robot & Mainly conceptual \\
SkySim~\cite{skysim2026} & Waypoints from the planner & Potential-field layer at 20~Hz; speed saturated at 0.5~m/s & Remote planner above the control layer; geofence in the planner & Gazebo with ROS~2; 3, 10 and 30 drones \\
Iannoli et al.~\cite{iannoli2026} & --- & --- & General-purpose models with planning tools and runtime guardrails & ArduPilot software-in-the-loop; four missions \\
Command\-Swarm~\cite{commandswarm2026} & Behaviour trees of whitelisted primitives & --- & Generates the tree; parser and safety classifier gate it & Trees scored, not executed \\
\midrule
This work & Formation slots with a \gls{pid} law; low-weighted global flocking terms & Potential field; geometric clamp at the integrator & One instance on the operator's device chooses a formation-level command; not in the control loop & Kinematic and physics simulators; five vehicles at 50~Hz \\
\bottomrule
\end{tabularx}
\end{table}

\paragraph{Critical comparison.} Table~\ref{tab:soa-swarm} separates the works on three axes. On
coordination, the foundational model produces an emergent flock with no assigned positions, while
the language-interface systems command explicit goals: waypoints in SkySim and behaviour trees of
whitelisted primitives in CommandSwarm. On separation, SkySim is the one language-interface system
that states its mechanism. It is a potential field, the method whose failure modes the analysis of
Koren and Borenstein catalogues, including the absence of a passage between closely spaced
obstacles. No prior row pairs the field with a mechanism that bounds what it fails to prevent. On
the place of the language model, SkySim falls outside both of Strobel et al.'s categories. Its model
neither synthesises the controller nor runs on each robot; it produces waypoints for a conventional
layer beneath it. The layering of a
language-model planner above a separate safety layer is therefore published. On evaluation, SkySim
and Iannoli et al.\ exercise their control layers in software-in-the-loop simulation, while
CommandSwarm does not execute its output. SkySim's planning latency of tens of seconds, judged
unsuitable for time-critical missions by its own authors, shows what the separation of layers leaves
unsolved: a layered design that routes every command through the planner stops only as fast as the
planner answers.

\section{Offline speech components}
\label{sec:speech-components}
The speech front ends of the systems reviewed above are of four kinds: recognisers run by a
platform or in the cloud, recognisers run offline in a cascade, direct recognisers that decide
without a transcript, and gates that react to voice activity alone. No comparative study was found
that selects among open components of these kinds on a four-core single-board processor under a
latency budget. For two of the components this work deploys, openWakeWord~\cite{oww} and
Silero~\gls{vad}~\cite{silero}, no peer-reviewed description was found at all; both are software
releases and are cited as such.

\paragraph{Recognisers in the reviewed systems.} The fixed-language systems use a recogniser as a
service: the Nuance platform for Landau and van Delden~\cite{landau2017}, and a cloud recogniser
corrected against domain phonemes for Contreras et al.~\cite{contreras2020}. Sikorski et al.\
recognise speech offline with VOSK~\cite{sikorski2025}. CommandSwarm chooses SeamlessM4T v2-large as
its speech front end, at about 4.0~s per utterance on unstated hardware. Whisper-medium speech
translation, at about 5.2~s, was the rejected alternative~\cite{commandswarm2026}. The cascade of
Henry et al.\ transcribes with Whisper before an intent classifier, and their direct recogniser
and the classifier of Sim\~oes et al.\ decide without a transcript~\cite{henry2026,simoes2024}.
MIRA's interruption gate reacts to voice activity alone, without transcription~\cite{mira}.

\paragraph{Open components.} Whisper is trained on 680{,}000 hours of weakly supervised audio, and
its smallest English-only model, \texttt{tiny.en}, has 39 million parameters~\cite{whisper}.
\texttt{whisper.cpp}~\cite{whispercpp} runs Whisper models in C and C++ on the \gls{cpu}, with no
Python runtime on the inference path. openWakeWord supplies a frozen feature front end that turns
audio into one speech embedding per 80~ms frame, over which small classifier heads can be trained
for arbitrary phrases~\cite{oww}. Silero \gls{vad} is a pre-trained neural voice-activity detector
distributed as an ONNX model that processes 32~ms windows with a recurrent state~\cite{silero}.
Table~\ref{tab:soa-speech} sets the front ends of the reviewed systems side by side.

\begin{table}[htbp]
\centering
\footnotesize
\caption[Speech front ends of spoken-command systems]{Speech front ends of the spoken-command
systems reviewed in this chapter. \emph{Offline}: whether recognition runs without a network
connection, as the source states it. \emph{Transcription-free path}: a path that decides from audio
without producing a transcript. \emph{Speech latency}: the time the source reports for its speech
stage or its whole pipeline. A dash marks a property this review did not establish from the source.
The last row is this work.}
\label{tab:soa-speech}
\begin{tabularx}{\textwidth}{@{}>{\raggedright\arraybackslash}p{2.3cm} >{\raggedright\arraybackslash}X >{\raggedright\arraybackslash}p{1.4cm} >{\raggedright\arraybackslash}p{2.4cm} >{\raggedright\arraybackslash}X@{}}
\toprule
System & Recogniser & Offline & Transcription-free path & Speech latency and hardware \\
\midrule
Landau and van Delden~\cite{landau2017} & Nuance platform & --- & No & --- \\
Contreras et al.~\cite{contreras2020} & Cloud recogniser, corrected against domain phonemes & No & No & --- \\
Sim\~oes et al.~\cite{simoes2024} & Recognition then a language model & --- & Direct classifier & Whole pipeline 1.233~s (cascade), 0.021~s (direct); --- \\
Henry et al.~\cite{henry2026} & Whisper, then CamemBERT-Large & --- & Direct intent recogniser & 202~ms (cascade); 7~ms on a \gls{gpu}, 106~ms on a \gls{cpu} (direct) \\
Sikorski et al.~\cite{sikorski2025} & VOSK & Yes & No & --- \\
Command\-Swarm~\cite{commandswarm2026} & SeamlessM4T v2-large; Whisper-medium rejected & Not claimed & No & About 4.0~s and 5.2~s per utterance; unstated hardware \\
MIRA~\cite{mira} & Not named & --- & Gate on voice activity & --- \\
\midrule
This work & \texttt{whisper.cpp tiny.en}, after a Silero \gls{vad} endpointer & Yes & Two-class keyword spotter on openWakeWord embeddings & Per stage, on the Raspberry~Pi~5 \gls{cpu} (Chapter~\ref{chap:validation}) \\
\bottomrule
\end{tabularx}
\end{table}

\paragraph{Critical comparison.} Table~\ref{tab:soa-speech} separates the front ends on three axes.
On where recognition runs, the older systems depend on a platform or a cloud service, and only
Sikorski et al.\ state offline recognition. CommandSwarm, the most recent system that accepts
speech, does not claim it. On a transcription-free path, three systems have one, and each uses it
differently. The two drone studies measure it as an alternative to the cascade, and MIRA uses voice
activity only to interrupt. None runs a direct recogniser beside a recogniser that transcribes. On
latency, the reported figures span nearly three orders of magnitude, from 7~ms for a direct recogniser on
a \gls{gpu} to about 4~s for a large multilingual model on unstated hardware. Only Henry et al.\
report the same component on two processors, and the move from \gls{gpu} to \gls{cpu} multiplies its
latency by about fifteen. A speech latency is therefore a property of the component and of its host
together, and most rows state one without the other. No prior row reports the latency of each
speech stage on the processor it deploys on, and none states a latency budget for its speech
stage.

\clearpage
\section{Research gaps}
\label{sec:research-gaps}

\paragraph{Synthesis.} Each of the four literatures reviewed in this chapter answers its own
question. Language interfaces to robots have moved from fixed command languages to language models
whose output is checked, and deterministic, automated checks appear in recent swarm
systems~\cite{iannoli2026,commandswarm2026}. Transcription-free paths have been measured to
cut the latency of a spoken command by one to two orders of magnitude at no cost in
accuracy~\cite{simoes2024,henry2026}. Fast paths have run concurrently with slow ones, gated by a
verifier or by preemption~\cite{relays2s,mira}. Language-model planners have been placed above a
potential-field layer that keeps vehicles apart~\cite{skysim2026}. Speech has been recognised offline, and a
quantised language model run offline on a control computer~\cite{sikorski2025}. The tables leave three gaps, stated below in the order the
design of Chapter~\ref{chap:architecture} closes them, after the published systems that already
cover parts of this work are conceded by name.

\paragraph{Conceded prior art.} Three systems cover parts of this work, and none of those parts is
claimed here. CommandSwarm covers spoken commands to a group of vehicles with the model's output
checked deterministically before it takes effect~\cite{commandswarm2026}. Its authors conclude that
parser acceptance and safety filtering remain necessary execution gates and that generation quality
alone is not sufficient, the position of Section~\ref{sec:safety-problem} reached independently.
SkySim covers the separation of a planner from a lower layer that keeps vehicles apart at the
control rate, and its fallback to holding position when a waypoint is rejected is close to the rule
of Section~\ref{sec:safety-problem} that a rejection becomes a hold~\cite{skysim2026}. MIRA covers an
uninterpreted stop followed by a hold, preemption of the slower path, and a control-rate safety layer
corresponding to the separation clamp at the integrator of this system~\cite{mira}. Constraining what
the operator may say is likewise established practice~\cite{landau2017,contreras2020}, and validation
downstream of an untrusted model is the arrangement of Silva and Burke~\cite{silva2026}. Like SkySim, and
outside both categories of Strobel et al.~\cite{llm2swarm2024}, this system uses one model
instance, on the operator's device, to map an utterance to a formation-level command for a
conventional controller.

\paragraph{Gap G1, what a fast path may carry.} The fast paths of Table~\ref{tab:soa-dual-path}
either have one action, as MIRA's abort does, or serve responsiveness or computation rather than
safety. The drone studies that measure a direct path choosing among commands build it as an
alternative to the cascade, never beside it. None of them states which commands a fast path may
carry. A fast path that chooses between commands without interpreting the utterance raises that
question, because a false accept on it publishes a well-formed command that no later check can tell
from a spoken one (Section~\ref{sec:safety-problem}). The second clause of the contribution of
Section~\ref{sec:contributions} closes this gap. The membership rule of
Section~\ref{sec:membership-rule} derives the fast path's commands from the cost of a false accept
and the value of the time saved. For this vocabulary it admits exactly two, a hold and an abort,
issued to five vehicles, and it states the condition under which a third could be admitted.

\paragraph{Gap G2, a stop with a budget of its own.} None of the systems in
Table~\ref{tab:voice-uav-survey} reports a stopping command handled separately from the path that
interprets commands. The layered design nearest to this one routes every command through a planner
whose latency runs to tens of seconds~\cite{skysim2026}. TypeFly answers the latency of generation
by shortening the path every command takes~\cite{typefly2024}. MIRA's gate stops the robot's motion
and preempts the slower path, but this review established no latency for it and no budget it is
held to. The first clause of the contribution closes
this gap: a reflex path with a latency budget of its own, 150~ms at p95 from the keyword offset,
which preempts language-model inference and must be met while the model decodes. The dual-path
decomposition (Section~\ref{sec:dual-path}) and the latency budget
(Section~\ref{sec:latency-budget}) specify it, and Section~\ref{sec:latency-experiment} measures it
idle and under load.

\paragraph{Gap G3, the whole chain on one board.} No system in Table~\ref{tab:voice-uav-survey}
executes its full chain on an \gls{sbc} without a \gls{gpu}, and no row of
Table~\ref{tab:soa-speech} reports per-stage speech latency on its deployment processor. The
offline arm of Sikorski et al.\ runs on a control computer, and its model exhibits significant
limitations in consistency and reliability~\cite{sikorski2025}; the \emph{M\'emoire de Master}
measures that cost for the models that can run on this device. The edge tier of Torkamani and Zarin
presumes that the network is available at least part of the time~\cite{torkamani2025}, which the
deployment site of Section~\ref{sec:operational-context} does not assure. Running the chain on one
board is engineering rather than a research contribution (Section~\ref{sec:contributions}), and an
architecture decision closes this gap: the allocation of the four cores of
Section~\ref{sec:resource-allocation}. It belongs among the gaps because it is the condition under
which the second must be closed. The reflex budget is met, or missed, on the same four cores on
which the language model decodes.

\paragraph{The baseline.} The agent of Lim et al.~\cite{lim2025} is the baseline for this document
because it is the published system nearest to this one in the part of the design the \emph{M\'emoire
de Master} fixes. It uses open language models of the same families, Qwen2.5 and Llama-3.2 among
them, hosted locally rather than by a cloud provider, and turns natural language into commands for a
PX4-based drone. The agent prompts its models zero-shot through Ollama, served from an RTX~3080~Ti
workstation on the local network, with no fine-tuning. Output is constrained by a sampling
temperature of 0.2 and by a prompt that fixes the format of turn and move commands. The mission
starts from a text message on ROS~2, and a single quadcopter is flown to a goal, with a
vision-language model detecting the target. Three of the four models produce valid flight commands
in all cases and DeepSeek-LLM in 38\%, the invalid outputs consisting of extraneous text, symbols and
markdown. The highest task success, 40\%, is obtained with Gemma3 and the vision-language model. The
fallback is a human operator who switches the vehicle to position mode, in which it hovers. Latency
is not reported.

\begin{table}[htbp]
\centering
\footnotesize
\caption[Design points of Lim et al.\ and of this work]{Design points of the agent of Lim et
al.~\cite{lim2025} and of this work. Every row differs, and each difference changes either the
latency or the accuracy a system can reach, which is why the two are compared by design and not by
measurement.}
\label{tab:lim-positioning}
\begin{tabularx}{\textwidth}{@{}>{\raggedright\arraybackslash}p{2.6cm} >{\raggedright\arraybackslash}X >{\raggedright\arraybackslash}X@{}}
\toprule
Property & Lim et al. & This work \\
\midrule
Input & Text message on ROS~2; no speech & Speech, recorded through the deployed microphone \\
Language model & Gemma3, Qwen2.5, Llama-3.2, DeepSeek-LLM & Qwen2.5-0.5B, quantised to Q4\_K\_M \\
Adaptation & Zero-shot prompting & Fine-tuned (\emph{M\'emoire de Master}) \\
Output constraint & Temperature 0.2; prompt fixing the command format & Decoding grammar; semantic validator; flight state machine \\
Inference hardware & RTX~3080~Ti workstation on the local network & Raspberry~Pi~5, four \gls{cpu} cores \\
Vehicles and task & One quadcopter flown to a goal & Five simulated vehicles commanded as a formation \\
Perception & Vision-language model detects the target & None \\
Stop & Human operator switches to position mode & Two-class keyword spotter publishes a hold or an abort \\
Reported & Valid commands; task success; detection rate & Per-stage latency; recognition under noise; safe-failure rate; formation accuracy; collisions \\
\bottomrule
\end{tabularx}
\end{table}

\paragraph{Comparison by design.} Table~\ref{tab:lim-positioning} sets the two systems side by
side, and they differ on every axis that determines latency or accuracy. Without speech input, the
figures of Lim et al.\ contain no recognition error and no endpointing wait, so neither the
acoustic-robustness experiment nor the parse-path latency of this document has a counterpart there.
A valid-command rate under zero-shot prompting measures how reliably a model follows a format it has
been asked to follow. Under a decoding grammar, validity is a property of the grammar, and the
\emph{M\'emoire de Master}, which owns that mechanism, reports it for the deployed parser. Any
latency, had one been reported, would be dominated by the difference between a desktop \gls{gpu}
and four \gls{cpu} cores. Task success is defined over goal-directed flight with visual detection, a
task this system does not perform, while formation accuracy is defined over a task Lim et al.\ do
not attempt. A stop that waits for a human operator offers nothing against which the reflex latency
of this system could be compared. A number from Lim et al.\ placed beside one from this document
would therefore measure the difference between two design points rather than a difference in
quality, and no such comparison is made. What the comparison establishes is the direction of each
difference: this system moves every axis of Table~\ref{tab:lim-positioning} towards the deployment of
Section~\ref{sec:operational-context}.

The gap this document fills is therefore narrower than any one of its parts suggests. Among the
systems reviewed, none combines deterministic checking of spoken swarm commands, a separate
control-rate safety layer, a stopping path that preempts the language model under a derived
membership rule and a budget of its own, and offline execution on an \gls{sbc}. The first two are
conceded to CommandSwarm and SkySim; the third is claimed against the fast paths of
Table~\ref{tab:soa-dual-path}; the fourth is the setting that turns the third into a problem of
processor allocation. Chapter~\ref{chap:architecture} designs the system that closes the three gaps,
and Chapter~\ref{chap:validation} measures it.
